
\begin{prop}[Deret Neumann]\label{prop_Neumann_series} Misalkan $a$ suatu elemen dari aljabar Banach beridentitas.
 \index{deret Neumann}%
 \index{deret!Neumann}%
Jika $\norm a < 1$, maka
$\vc 1 - a \in \inv A$ dan $(\vc 1 - a)^{-1} = \sum_{k=0}^\infty a^k$.
\end{prop}

\begin{proof}[\emph{Petunjuk untuk bukti}] Tunjukkan bahwa barisan $(\vc 1, a, a^2, \dots)$
dapat dijumlahkan. (Dalam aljabar beridentitas kita mengartikan $a^0$ sebagai~$\vc 1$.)  \ns
\end{proof}

\begin{cor}\label{cor_Neumann_series} Jika $a$ suatu elemen dari aljabar Banach beridentitas dengan $\norm{\vc 1 - a} < 1$, maka $a$ invertibel dan
    \[ \norm{a^{-1}} \le \frac{1}{1 - \norm{\vc 1 - a}}\,.  \]
\end{cor}

\begin{cor}\label{cor2_Neumann_series} Misalkan $a$ suatu elemen invertibel dari aljabar Banach beridentitas $A$ dan $b \in A$,
serta $\norm{a - b} < \norm{a^{-1}}^{-1}$. Maka $b$ invertibel dalam~$A$.
\end{cor}

\begin{proof}[\emph{Petunjuk untuk bukti}] Gunakan korolari sebelumnya untuk menunjukkan bahwa $a^{-1}b$ invertibel.
\ns \end{proof}

\begin{cor} Jika $a$ termasuk dalam suatu aljabar Banach beridentitas dan $\norm a < 1$, maka
   \[ \norm{(\vc 1 - a)^{-1} - \vc 1} \le \frac{\norm a}{1 - \norm a}\,. \]
\end{cor}

\begin{prop}\label{000644fa} Jika $A$ aljabar Banach beridentitas, maka $\open{\inv A}A$.
\end{prop}

\begin{proof}[\emph{Petunjuk untuk bukti}] Misalkan $a \in \inv A$. Tunjukkan, untuk $h$ yang cukup kecil,
bahwa $\vc 1 - a^{-1}h$ invertibel.  \ns
\end{proof}

\begin{prop}\label{000646fa} Misalkan $A$ aljabar Banach beridentitas. Pemetaan $a \mapsto a^{-1}$ dari $\inv A$ ke
dirinya sendiri merupakan homeomorfisme.
\end{prop}

\begin{prop} Misalkan $A$ aljabar Banach beridentitas. Pemetaan $r\colon a \mapsto a^{-1}$ dari
$\inv A$ ke dirinya sendiri bersifat diferensiabel dan pada setiap elemen invertibel $a$, berlaku
$dr_a(h) = -a^{-1}ha^{-1}$ untuk setiap $h \in A$.
\end{prop}

\begin{proof}[\emph{Petunjuk untuk bukti}] Untuk pengantar singkat mengenai diferensiasi pada ruang
berdimensi tak hingga, lihat Bab 13 dari catatan saya~\cite{Erdman:2007} tentang Analisis Real.  \ns
\end{proof}

\begin{prop}\label{000661} Misalkan $a$ suatu elemen dari aljabar Banach beridentitas~$A$. Maka spektrum $a$
kompak dan $\abs\lambda \le \norm a$ untuk setiap $\lambda \in \sigma(a)$.
\end{prop}

\begin{proof}[\emph{Petunjuk untuk bukti}] Gunakan \emph{teorema Heine-Borel}. Untuk membuktikan bahwa
spektrum tertutup, perhatikan bahwa $(\sigma(a))^c = f^\gets(\inv A)$ dengan $f(\lambda) = a -
\lambda\mathbf 1$ untuk setiap bilangan kompleks~$\lambda$. Tunjukkan juga bahwa jika $\abs\lambda >
\norm a$, maka $\mathbf 1 - \lambda^{-1}a$ invertibel.  \ns
\end{proof}

\begin{defn} Misalkan $a$ suatu elemen dari aljabar Banach beridentitas.
 \index{resolven!pemetaan}%
\df{Pemetaan resolven} untuk $a$ didefinisikan oleh
   \[ R_a \colon \C \setminus \sigma(a) \sto A
                 \colon \lambda \mapsto (a  - \lambda\vc 1)^{-1}\,. \]
\end{defn}

\begin{defn} Misalkan $\open U\C$ dan $A$ suatu aljabar Banach beridentitas. Suatu fungsi $f \colon U \sto A$
 \index{analitik}%
bersifat \df{analitik} pada $U$ jika
   \[ f'(a) := \lim_{z \sto a} \frac{f(z) - f(a)}{z - a} \]
ada untuk setiap $a \in U$. Suatu fungsi bernilai kompleks yang analitik pada seluruh $\C$
disebut fungsi
 \index{entire}
\df{entire}.
\end{defn}

\begin{prop} Untuk $a$ suatu elemen dari aljabar Banach beridentitas $A$ dan $\phi$ suatu fungsional linear
terbatas pada~$A$, misalkan $f := \phi \circ R_a \colon \C \setminus \sigma(a) \sto \C$. Maka
 \begin{enumerate}[label=\rm{(\alph*)}]
  \item $f$ analitik pada domainnya; dan
  \item $f(\lambda) \sto 0$ ketika $\abs\lambda \sto \infty$.
 \end{enumerate}
\end{prop}

\begin{proof}[\emph{Petunjuk untuk bukti}] Untuk (i), perhatikan bahwa dalam aljabar beridentitas $a^{-1} - b^{-1}
= a^{-1}(b - a)b^{-1}$.   \ns
\end{proof}

\begin{thm}[Liouville]\label{C073131} Setiap fungsi entire yang terbatas pada $\C$ bersifat konstan.
 \index{teorema Liouville}%
\end{thm}

\begin{proof}[\emph{Komentar tentang bukti}] Pembuktian teorema ini memerlukan sedikit latar belakang
dalam teori fungsi analitik, yakni materi yang tidak dibahas dalam catatan ini. Teorema
tersebut tentu dibahas dalam setiap mata kuliah pengantar tentang variabel kompleks. Sebagai contoh, bukti yang
baik dapat ditemukan dalam \cite{Rudin:1987}, teorema 10.23.   \ns
\end{proof}

\begin{prop}\label{C073134} Spektrum setiap elemen dari aljabar Banach beridentitas tidak kosong.
\end{prop}

\begin{proof}[\emph{Petunjuk untuk bukti}] Gunakan argumen kontradiksi. Gunakan \emph{teorema Liouville}~\ref{C073131}
untuk menunjukkan bahwa $\phi \circ R_a$ konstan bagi setiap fungsional linear terbatas $\phi$ pada~$A$. Kemudian gunakan
(salah satu versi) \emph{teorema Hahn-Banach} untuk membuktikan bahwa $R_a$ konstan. Mengapa konstanta ini harus~$0$?   \ns
\end{proof}
